\Needspace{8\baselineskip}
\subsection{1.1 Impulse seed theorem}
\label{result:1.1}
\textbf{Traditional locator.} Theorem 1.1, final paper p.~8; \code{result:01}.
\textbf{Exact reviewed scope.} Complex l2 impulse expansion, finite-support identity, unitary shifts and dense shift-generated span.
{\small\code{KomorebiPaperV.HilbertPackets.impulse_expansion}\par}
{\footnotesize\color{muted}\code{KomorebiPaperV/HilbertPackets.lean}, line 135.\par}\smallskip
{\small\code{KomorebiPaperV.HilbertPackets.finite_impulse_expansion}\par}
{\footnotesize\color{muted}\code{KomorebiPaperV/HilbertPackets.lean}, line 207.\par}\smallskip
{\small\code{KomorebiPaperV.HilbertPackets.shift_generated_dense_span}\par}
{\footnotesize\color{muted}\code{KomorebiPaperV/HilbertPackets.lean}, line 203.\par}\smallskip
{\small\code{KomorebiPaperV.HilbertPackets.shift_impulse}\par}
{\footnotesize\color{muted}\code{KomorebiPaperV/HilbertPackets.lean}, line 183.\par}\smallskip
\Needspace{8\baselineskip}
\subsection{2.1 Polyphase projector theorem}
\label{result:2.1}
\textbf{Traditional locator.} Theorem 2.1, final paper p.~9; \code{result:02}.
\textbf{Exact reviewed scope.} Bounded orthogonal residue projectors on complex l2.
{\small\code{KomorebiPaperV.HilbertPackets.projector_idempotent}\par}
{\footnotesize\color{muted}\code{KomorebiPaperV/HilbertPackets.lean}, line 56.\par}\smallskip
{\small\code{KomorebiPaperV.HilbertPackets.projector_annihilate}\par}
{\footnotesize\color{muted}\code{KomorebiPaperV/HilbertPackets.lean}, line 62.\par}\smallskip
{\small\code{KomorebiPaperV.HilbertPackets.projector_complete}\par}
{\footnotesize\color{muted}\code{KomorebiPaperV/HilbertPackets.lean}, line 69.\par}\smallskip
{\small\code{KomorebiPaperV.HilbertPackets.projector_symmetric}\par}
{\footnotesize\color{muted}\code{KomorebiPaperV/HilbertPackets.lean}, line 139.\par}\smallskip
{\small\code{KomorebiPaperV.HilbertPackets.projector_orthogonal}\par}
{\footnotesize\color{muted}\code{KomorebiPaperV/HilbertPackets.lean}, line 147.\par}\smallskip
\Needspace{8\baselineskip}
\subsection{2.2 Analysis–synthesis identity}
\label{result:2.2}
\textbf{Traditional locator.} Corollary 2.2, final paper p.~9; \code{result:03}.
\textbf{Exact reviewed scope.} Signed integer polyphase linear isometry; channels represented by joint l2 coordinates.
{\small\code{KomorebiPaperV.HilbertPackets.polyphase_synthesis_analysis}\par}
{\footnotesize\color{muted}\code{KomorebiPaperV/HilbertPackets.lean}, line 123.\par}\smallskip
{\small\code{KomorebiPaperV.HilbertPackets.polyphase_analysis_synthesis}\par}
{\footnotesize\color{muted}\code{KomorebiPaperV/HilbertPackets.lean}, line 127.\par}\smallskip
{\small\code{KomorebiPaperV.HilbertPackets.polyphase_energy}\par}
{\footnotesize\color{muted}\code{KomorebiPaperV/HilbertPackets.lean}, line 131.\par}\smallskip
\Needspace{8\baselineskip}
\subsection{2.3 Exact nested refinement}
\label{result:2.3}
\textbf{Traditional locator.} Theorem 2.3, final paper p.~10; \code{result:04}.
\textbf{Exact reviewed scope.} Literal divisible-modulus bounded projector refinement.
{\small\code{KomorebiPaperV.HilbertPackets.residue_projector_refinement}\par}
{\footnotesize\color{muted}\code{KomorebiPaperV/HilbertPackets.lean}, line 187.\par}\smallskip
\Needspace{8\baselineskip}
\subsection{3.1 Exact FIR polyphase transfer law}
\label{result:3.1}
\textbf{Traditional locator.} Theorem 3.1, final paper p.~11; \code{result:05}.
\textbf{Exact reviewed scope.} Time-domain transfer for every bilateral sequence; channel z-transform identity for z nonzero wherever each weighted input series is summable.
{\small\code{KomorebiPaperV.ZTransform.fir_channel_transform}\par}
{\footnotesize\color{muted}\code{KomorebiPaperV/ZTransform.lean}, line 49.\par}\smallskip
{\small\code{KomorebiPaperV.ZTransform.fir_matrix_transform}\par}
{\footnotesize\color{muted}\code{KomorebiPaperV/ZTransform.lean}, line 64.\par}\smallskip
{\small\code{KomorebiPaperV.FIRTransfer.fir_polyphase_transfer}\par}
{\footnotesize\color{muted}\code{KomorebiPaperV/FIRTransfer.lean}, line 96.\par}\smallskip
\Needspace{8\baselineskip}
\subsection{4.1 Perfect-reconstruction criterion}
\label{result:4.1}
\textbf{Traditional locator.} Theorem 4.1, final paper p.~13; \code{result:06}.
\textbf{Exact reviewed scope.} Concrete finite Laurent/FIR banks on every bilateral sequence, faithful matrix/operator correspondence, gain and channel delay; original-sample delay is N*d. Arbitrary scalar-matrix criterion is also retained in FiniteTransforms.
{\small\code{KomorebiPaperV.FilterBank.act_is_fir}\par}
{\footnotesize\color{muted}\code{KomorebiPaperV/FilterBank.lean}, line 36.\par}\smallskip
{\small\code{KomorebiPaperV.FilterBank.bank_mul}\par}
{\footnotesize\color{muted}\code{KomorebiPaperV/FilterBank.lean}, line 84.\par}\smallskip
{\small\code{KomorebiPaperV.FilterBank.bank_faithful}\par}
{\footnotesize\color{muted}\code{KomorebiPaperV/FilterBank.lean}, line 104.\par}\smallskip
{\small\code{KomorebiPaperV.FilterBank.perfect_reconstruction_iff}\par}
{\footnotesize\color{muted}\code{KomorebiPaperV/FilterBank.lean}, line 133.\par}\smallskip
{\small\code{KomorebiPaperV.FilterBank.bank_transform}\par}
{\footnotesize\color{muted}\code{KomorebiPaperV/FilterBank.lean}, line 147.\par}\smallskip
{\small\code{KomorebiPaperV.FilterBank.original_sample_delay}\par}
{\footnotesize\color{muted}\code{KomorebiPaperV/FilterBank.lean}, line 155.\par}\smallskip
\Needspace{8\baselineskip}
\subsection{4.2 Haar energy identity}
\label{result:4.2}
\textbf{Traditional locator.} Corollary 4.2, final paper p.~14; \code{result:07}.
\textbf{Exact reviewed scope.} Literal interleaved input and normalized low/detail channels in complex l2.
{\small\code{KomorebiPaperV.HaarChannels.haar_energy}\par}
{\footnotesize\color{muted}\code{KomorebiPaperV/HaarChannels.lean}, line 40.\par}\smallskip
{\small\code{KomorebiPaperV.HaarChannels.haar_low_samples}\par}
{\footnotesize\color{muted}\code{KomorebiPaperV/HaarChannels.lean}, line 44.\par}\smallskip
{\small\code{KomorebiPaperV.HaarChannels.haar_high_samples}\par}
{\footnotesize\color{muted}\code{KomorebiPaperV/HaarChannels.lean}, line 48.\par}\smallskip
\Needspace{8\baselineskip}
\subsection{4.3 Exact lifting reconstruction}
\label{result:4.3}
\textbf{Traditional locator.} Theorem 4.3, final paper p.~15; \code{result:08}.
\textbf{Exact reviewed scope.} Pairwise predict/update inverse and determinant one, with both flat finite-signal inverse laws.
{\small\code{KomorebiPaperV.FiniteTransforms.unlift_lift_flat}\par}
{\footnotesize\color{muted}\code{KomorebiPaperV/FiniteTransforms.lean}, line 78.\par}\smallskip
{\small\code{KomorebiPaperV.FiniteTransforms.lift_unlift_flat}\par}
{\footnotesize\color{muted}\code{KomorebiPaperV/FiniteTransforms.lean}, line 83.\par}\smallskip
{\small\code{KomorebiPaperV.lifting_synthesis_analysis}\par}
{\footnotesize\color{muted}\code{KomorebiPaperV/Lifting.lean}, line 21.\par}\smallskip
{\small\code{KomorebiPaperV.lifting_analysis_synthesis}\par}
{\footnotesize\color{muted}\code{KomorebiPaperV/Lifting.lean}, line 25.\par}\smallskip
{\small\code{KomorebiPaperV.lifting_determinant}\par}
{\footnotesize\color{muted}\code{KomorebiPaperV/Lifting.lean}, line 29.\par}\smallskip
\Needspace{8\baselineskip}
\subsection{5.1 Wavelet-packet reconstruction and energy}
\label{result:5.1}
\textbf{Traditional locator.} Theorem 5.1, final paper p.~15; \code{result:09}.
\textbf{Exact reviewed scope.} All finite depths of literal Haar packet operators; sums of true Hilbert adjoints reconstruct, and packet norms preserve energy.
{\small\code{KomorebiPaperV.WaveletPackets.packet_energy}\par}
{\footnotesize\color{muted}\code{KomorebiPaperV/WaveletPackets.lean}, line 68.\par}\smallskip
{\small\code{KomorebiPaperV.WaveletPackets.packet_operator_reconstruction}\par}
{\footnotesize\color{muted}\code{KomorebiPaperV/WaveletPackets.lean}, line 94.\par}\smallskip
{\small\code{KomorebiPaperV.WaveletPackets.packet_reconstruction}\par}
{\footnotesize\color{muted}\code{KomorebiPaperV/WaveletPackets.lean}, line 98.\par}\smallskip
\Needspace{8\baselineskip}
\subsection{6.1 Exact known-edge theorem}
\label{result:6.1}
\textbf{Traditional locator.} Theorem 6.1, final paper p.~16; \code{result:10}.
\textbf{Exact reviewed scope.} Exact known-edge cancellation in arbitrary seminormed additive groups.
{\small\code{KomorebiPaperV.known_edge_error}\par}
{\footnotesize\color{muted}\code{KomorebiPaperV/KnownEdge.lean}, line 17.\par}\smallskip
{\small\code{KomorebiPaperV.known_edge_norm}\par}
{\footnotesize\color{muted}\code{KomorebiPaperV/KnownEdge.lean}, line 34.\par}\smallskip
\Needspace{8\baselineskip}
\subsection{6.2 Several exact edge packets}
\label{result:6.2}
\textbf{Traditional locator.} Corollary 6.2, final paper p.~16; \code{result:11}.
\textbf{Exact reviewed scope.} Finite families of exactly retained jump packets.
{\small\code{KomorebiPaperV.several_known_edges_error}\par}
{\footnotesize\color{muted}\code{KomorebiPaperV/KnownEdge.lean}, line 25.\par}\smallskip
{\small\code{KomorebiPaperV.several_known_edges_norm}\par}
{\footnotesize\color{muted}\code{KomorebiPaperV/KnownEdge.lean}, line 38.\par}\smallskip
\Needspace{8\baselineskip}
\subsection{7.1 Packet perturbation bound}
\label{result:7.1}
\textbf{Traditional locator.} Theorem 7.1, final paper p.~17; \code{result:12}.
\textbf{Exact reviewed scope.} Measurable scalar masks in arbitrary measure spaces; extended finite-p norms, and pointwise plus essential supremum bounds. Finite mask and symmetric-difference measures yield the finite real-valued statement.
{\small\code{KomorebiPaperV.MaskPerturbation.indicator_difference_norm}\par}
{\footnotesize\color{muted}\code{KomorebiPaperV/MaskPerturbation.lean}, line 25.\par}\smallskip
{\small\code{KomorebiPaperV.MaskPerturbation.finite_p_perturbation}\par}
{\footnotesize\color{muted}\code{KomorebiPaperV/MaskPerturbation.lean}, line 43.\par}\smallskip
{\small\code{KomorebiPaperV.MaskPerturbation.pointwise_sup_bound}\par}
{\footnotesize\color{muted}\code{KomorebiPaperV/MaskPerturbation.lean}, line 59.\par}\smallskip
{\small\code{KomorebiPaperV.MaskPerturbation.infinity_perturbation}\par}
{\footnotesize\color{muted}\code{KomorebiPaperV/MaskPerturbation.lean}, line 73.\par}\smallskip
{\small\code{KomorebiPaperV.MaskPerturbation.full_signal_perturbation}\par}
{\footnotesize\color{muted}\code{KomorebiPaperV/MaskPerturbation.lean}, line 79.\par}\smallskip
\Needspace{8\baselineskip}
\subsection{11.2 Exact synthesis from every complete pruning}
\label{result:11.2}
\textbf{Traditional locator.} Theorem 11.2, final paper p.~24; \code{result:13}.
\textbf{Exact reviewed scope.} Literal finite node sums over covering antichains, with exact operator and applied-signal synthesis under the manuscript conservation hypotheses.
{\small\code{KomorebiPaperV.PruningSums.cutSum_eq_nodeSum}\par}
{\footnotesize\color{muted}\code{KomorebiPaperV/PruningSums.lean}, line 54.\par}\smallskip
{\small\code{KomorebiPaperV.PruningSums.antichain_synthesis}\par}
{\footnotesize\color{muted}\code{KomorebiPaperV/PruningSums.lean}, line 81.\par}\smallskip
{\small\code{KomorebiPaperV.PruningSums.packet_projector_synthesis}\par}
{\footnotesize\color{muted}\code{KomorebiPaperV/PruningSums.lean}, line 88.\par}\smallskip
\Needspace{8\baselineskip}
\subsection{11.3 Finite globally optimal stable selector}
\label{result:11.3}
\textbf{Traditional locator.} Theorem 11.3, final paper p.~24; \code{result:14}.
\textbf{Exact reviewed scope.} Named nonadditive representation/leakage/complexity criterion; attained positive gap when competitors exist, uniform perturbation stability, and singleton case. Every pruning preserves synthesis.
{\small\code{KomorebiPaperV.PruningSums.objective_minimum}\par}
{\footnotesize\color{muted}\code{KomorebiPaperV/PruningSums.lean}, line 134.\par}\smallskip
{\small\code{KomorebiPaperV.PruningSums.objective_gap_attained}\par}
{\footnotesize\color{muted}\code{KomorebiPaperV/PruningSums.lean}, line 165.\par}\smallskip
{\small\code{KomorebiPaperV.PruningSums.objective_stability}\par}
{\footnotesize\color{muted}\code{KomorebiPaperV/PruningSums.lean}, line 148.\par}\smallskip
{\small\code{KomorebiPaperV.PruningSums.singleton_selector}\par}
{\footnotesize\color{muted}\code{KomorebiPaperV/PruningSums.lean}, line 158.\par}\smallskip
{\small\code{KomorebiPaperV.PruningSums.packet_projector_synthesis}\par}
{\footnotesize\color{muted}\code{KomorebiPaperV/PruningSums.lean}, line 88.\par}\smallskip
\Needspace{8\baselineskip}
\subsection{11.4 Dynamic-programming best packet basis}
\label{result:11.4}
\textbf{Traditional locator.} Theorem 11.4, final paper p.~25; \code{result:15}.
\textbf{Exact reviewed scope.} Cached bottom-up evaluator equals the global selector; stop on ties. One comparison per internal node; at most 3*node-count unit operations for visits/comparisons/value additions with node costs supplied. Not a Lean runtime or scalar-cost evaluation bound.
{\small\code{KomorebiPaperV.PruningAlgorithm.solve_correct}\par}
{\footnotesize\color{muted}\code{KomorebiPaperV/PruningAlgorithm.lean}, line 36.\par}\smallskip
{\small\code{KomorebiPaperV.PruningAlgorithm.solve_comparisons}\par}
{\footnotesize\color{muted}\code{KomorebiPaperV/PruningAlgorithm.lean}, line 52.\par}\smallskip
{\small\code{KomorebiPaperV.PruningAlgorithm.solve_visits}\par}
{\footnotesize\color{muted}\code{KomorebiPaperV/PruningAlgorithm.lean}, line 61.\par}\smallskip
{\small\code{KomorebiPaperV.PruningAlgorithm.solve_additions}\par}
{\footnotesize\color{muted}\code{KomorebiPaperV/PruningAlgorithm.lean}, line 69.\par}\smallskip
{\small\code{KomorebiPaperV.PruningAlgorithm.solve_linear_work}\par}
{\footnotesize\color{muted}\code{KomorebiPaperV/PruningAlgorithm.lean}, line 89.\par}\smallskip
{\small\code{KomorebiPaperV.PruningAlgorithm.computed_antichain_optimal}\par}
{\footnotesize\color{muted}\code{KomorebiPaperV/PruningAlgorithm.lean}, line 97.\par}\smallskip
\Needspace{8\baselineskip}
\subsection{11.5 Local perturbation certificate}
\label{result:11.5}
\textbf{Traditional locator.} Corollary 11.5, final paper p.~26; \code{result:16}.
\textbf{Exact reviewed scope.} Exact maximum antichain error envelope and strict-gap stability; probability transfer from a joint local-error event is proved in PruningAlgorithm.
{\small\code{KomorebiPaperV.PruningSums.antichain_perturbation}\par}
{\footnotesize\color{muted}\code{KomorebiPaperV/PruningSums.lean}, line 108.\par}\smallskip
{\small\code{KomorebiPaperV.PruningSums.antichain_envelope_attained}\par}
{\footnotesize\color{muted}\code{KomorebiPaperV/PruningSums.lean}, line 115.\par}\smallskip
{\small\code{KomorebiPaperV.PruningSums.antichain_selector_stable}\par}
{\footnotesize\color{muted}\code{KomorebiPaperV/PruningSums.lean}, line 120.\par}\smallskip
{\small\code{KomorebiPaperV.PruningAlgorithm.antichain_probability}\par}
{\footnotesize\color{muted}\code{KomorebiPaperV/PruningAlgorithm.lean}, line 102.\par}\smallskip
\Needspace{8\baselineskip}
\subsection{12.1 Cyclic product law}
\label{result:12.1}
\textbf{Traditional locator.} Theorem 12.1, final paper p.~28; \code{result:17}.
\textbf{Exact reviewed scope.} Formal power-series cyclic product and literal q-tuple power formula over commutative semirings, including q=0. Analytic interpretations apply on their convergence domains; no arbitrary divergent evaluation is asserted.
{\small\code{KomorebiPaperV.CyclicSeries.cyclic_product_law}\par}
{\footnotesize\color{muted}\code{KomorebiPaperV/CyclicSeries.lean}, line 126.\par}\smallskip
{\small\code{KomorebiPaperV.CyclicSeries.cyclic_power_law}\par}
{\footnotesize\color{muted}\code{KomorebiPaperV/CyclicSeries.lean}, line 195.\par}\smallskip
\Needspace{8\baselineskip}
\subsection{13.2 Polyphase Möbius action}
\label{result:13.2}
\textbf{Traditional locator.} Theorem 13.2, final paper p.~30; \code{result:18}.
\textbf{Exact reviewed scope.} Explicit exponential root of unity and modulo-N action on the projective quotient of nonzero homogeneous pairs; Cayley conjugacy and the finite-chart formula.
{\small\code{KomorebiPaperV.SmithCyclic.omega_pow}\par}
{\footnotesize\color{muted}\code{KomorebiPaperV/SmithCyclic.lean}, line 12.\par}\smallskip
{\small\code{KomorebiPaperV.SmithCyclic.cyclic_composition}\par}
{\footnotesize\color{muted}\code{KomorebiPaperV/SmithCyclic.lean}, line 78.\par}\smallskip
{\small\code{KomorebiPaperV.SmithCyclic.cayley_conjugacy}\par}
{\footnotesize\color{muted}\code{KomorebiPaperV/SmithCyclic.lean}, line 87.\par}\smallskip
{\small\code{KomorebiPaperV.SmithCyclic.cyclic_chart}\par}
{\footnotesize\color{muted}\code{KomorebiPaperV/SmithCyclic.lean}, line 92.\par}\smallskip
{\small\code{KomorebiPaperV.SmithCyclic.projective_mul}\par}
{\footnotesize\color{muted}\code{KomorebiPaperV/SmithCyclic.lean}, line 56.\par}\smallskip
\Needspace{8\baselineskip}
\subsection{13.3 Impedance–admittance is the two-phase split}
\label{result:13.3}
\textbf{Traditional locator.} Corollary 13.3, final paper p.~31; \code{result:19}.
\textbf{Exact reviewed scope.} The N=2 phase is -1; its action swaps homogeneous coordinates, exchanging zero/infinity and giving inverse impedance on the nonzero chart.
{\small\code{KomorebiPaperV.SmithCyclic.omega_two}\par}
{\footnotesize\color{muted}\code{KomorebiPaperV/SmithCyclic.lean}, line 19.\par}\smallskip
{\small\code{KomorebiPaperV.SmithCyclic.two_phase_swap}\par}
{\footnotesize\color{muted}\code{KomorebiPaperV/SmithCyclic.lean}, line 97.\par}\smallskip
{\small\code{KomorebiPaperV.SmithCyclic.two_phase_impedance_admittance}\par}
{\footnotesize\color{muted}\code{KomorebiPaperV/SmithCyclic.lean}, line 100.\par}\smallskip
\Needspace{8\baselineskip}
\subsection{14.1 Dyadic–odd packet factorization}
\label{result:14.1}
\textbf{Traditional locator.} Proposition 14.1, final paper p.~34; \code{result:20}.
\textbf{Exact reviewed scope.} Signed integer CRT residue masks and dyadic-odd factorization.
{\small\code{KomorebiPaperV.CRT.residue_mask_factorization}\par}
{\footnotesize\color{muted}\code{KomorebiPaperV/CRT.lean}, line 20.\par}\smallskip
{\small\code{KomorebiPaperV.CRT.dyadic_odd_packet_factorization}\par}
{\footnotesize\color{muted}\code{KomorebiPaperV/CRT.lean}, line 34.\par}\smallskip
